\documentclass[12pt, a4paper]{article}

% Pacotes essenciais
\usepackage[utf8]{inputenc}
\usepackage[T1]{fontenc}
\usepackage[brazil]{babel}
\usepackage[margin=2.5cm]{geometry}
\usepackage{booktabs}
\usepackage{setspace}

% Metadados do documento
\title{\vspace{-2cm}\textbf{Proposta de Projeto de Fim de Curso I}\\ \large Reavaliação do Algoritmo de Agrupamento \textit{microTEDAclus} e Proposta de um Modelo Aprimorado}
\author{\textbf{Aluno:} Eduardo Henrique Basilio de Carvalho \\ \textbf{Orientador:} Prof. Janier Arias Garcia}
\date{\today}

\begin{document}

\maketitle
\onehalfspacing

\section{Objetivo}
Investigar e expor formalmente as inconsistências teóricas e algorítmicas do modelo \textit{microTEDAclus}, propondo uma arquitetura alternativa fundamentada nos mesmos princípios estatísticos e matemáticos, porém com uma aplicação analítica corrigida e distinta.

\section{Justificativa}
A motivação para este estudo fundamenta-se na constatação empírica de divergências na literatura recente. Tanto a replicação da metodologia proposta no artigo \textit{Evolving clustering algorithm based on mixture of typicalities for stream data mining}, quanto a auditoria dos arquivos de implementação referentes ao trabalho \textit{Mineração de dados aplicada a problemas de last-mile na logística}, revelaram incompatibilidade estrita entre os resultados reportados pelos autores e o comportamento observado em experimentação independente. Faz-se necessária uma correção formal do modelo para garantir a confiabilidade de sua aplicação.

\section{Metodologia}
A execução do projeto ocorrerá por meio das seguintes etapas sistemáticas:
\begin{enumerate}
    \item \textbf{Revisão Bibliográfica:} Mapeamento do estado da arte e análise da fundamentação matemática pertinente aos modelos derivados do \textit{TEDAclus}.
    \item \textbf{Experimentos Sistemáticos nos Modelos Originais:} Condução de testes padronizados e testes de estresse nas implementações originais para isolamento e prova das falhas supracitadas.
    \item \textbf{Desenvolvimento do Modelo Alternativo:} Estruturação matemática e algorítmica da nova proposta.
    \item \textbf{Comparação Sistemática:} Avaliação de desempenho e estabilidade estatística, contrastando o modelo aprimorado com seus predecessores.
\end{enumerate}

\section{Cronograma de Entregas}
As metas estabelecidas para a primeira fase do projeto respeitarão o cronograma estipulado na Tabela \ref{tab:cronograma}.

\begin{table}[h!]
\centering
\caption{Cronograma de entregas da primeira fase do Projeto de Fim de Curso.}
\vspace{0.2cm}
\label{tab:cronograma}
\begin{tabular}{@{}ll@{}}
\toprule
\textbf{Data (2026)} & \textbf{Marco de Entrega / Artefato} \\ \midrule
06 de Abril & Relatório de Revisão Bibliográfica. \\
11 de Maio  & Relatório de Testes e Análise Crítica dos Modelos Originais. \\
08 de Junho & Descrição Formal e Modelagem Matemática do Modelo Alternativo. \\ \bottomrule
\end{tabular}
\end{table}

\end{document}
